
\documentclass[11pt,a4paper]{article}

\usepackage[margin=1in]{geometry}
\usepackage{amsmath,amssymb}
\usepackage{siunitx}
\usepackage{booktabs}
\usepackage{array}
\usepackage{enumitem}
\usepackage{hyperref}
\usepackage{fancyhdr}

\hypersetup{
    colorlinks=true,
    linkcolor=blue,
    citecolor=blue,
    urlcolor=blue
}

\pagestyle{fancy}
\fancyhf{}
\lhead{BPhO Computational Physics Challenge 2026}
\rhead{Quantum Mechanics 2}
\cfoot{\thepage}

\title{\textbf{Quantum Mechanics 2}\\Worked Solutions}
\author{BPhO Computational Physics Challenge 2026}
\date{}

\begin{document}

\maketitle

\section*{Introduction}

This document gives worked solutions to the \emph{Quantum Mechanics 2} problem sheet.
The calculations use the constants supplied on the sheet and have been checked against
the accompanying BPhO solution document. The aim is to show the physical reasoning,
mathematical derivations and numerical substitutions rather than simply state the final
answers.

% ============================================================
\section{Question 1 --- The Caesium-133 clock transition}

The definition of the SI second is based on radiation from Caesium-133 with frequency
\[
f=9\,192\,631\,770\ \mathrm{Hz}.
\]

\subsection{(i) Wavelength}

For electromagnetic radiation,
\[
c=f\lambda,
\]
so
\[
\lambda=\frac{c}{f}.
\]

Using
\[
c=2.998\times10^8\ \mathrm{m\,s^{-1}},
\]
we obtain
\[
\lambda
=
\frac{2.998\times10^8}
{9.192631770\times10^9}
\approx3.26\times10^{-2}\ \mathrm m.
\]

Thus
\[
\boxed{\lambda\approx3.26\ \mathrm{cm}}.
\]

A wavelength of a few centimetres corresponds to the
\[
\boxed{\text{microwave}}
\]
region of the electromagnetic spectrum.

\subsection{(ii) Energy change}

A photon has energy
\[
E=hf.
\]

The energy in electron-volts is
\[
E_{\rm eV}=\frac{hf}{e}.
\]

Therefore
\[
E_{\rm eV}
=
\frac{(6.63\times10^{-34})(9.192631770\times10^9)}
{1.60\times10^{-19}}
\approx3.81\times10^{-5}\ \mathrm{eV}.
\]

Hence
\[
\boxed{\Delta E\approx3.81\times10^{-5}\ \mathrm{eV}}.
\]

This is an \emph{electronic} rather than a nuclear energy scale. The relevant
transition is between very closely spaced atomic states of Caesium; nuclear
transitions typically involve energies of order keV--MeV, vastly larger than this
microwave transition.

% ============================================================
\section{Question 2 --- Photoelectric effect}

The photoelectric equation is
\[
K_{\max}=hf-\phi,
\]
where $\phi$ is the work function.

\subsection{(i) Work function}

No electrons are emitted for wavelengths greater than \SI{300}{nm}. Thus the
threshold wavelength is
\[
\lambda_0=\SI{300}{nm}.
\]

At threshold the emitted electron has zero kinetic energy:
\[
hf_0=\phi.
\]

Since
\[
f_0=\frac{c}{\lambda_0},
\]
we obtain
\[
\phi=\frac{hc}{\lambda_0}.
\]

Thus
\[
\phi
=
\frac{(6.63\times10^{-34})(2.998\times10^8)}
{300\times10^{-9}}
\approx6.63\times10^{-19}\ \mathrm J.
\]

Converting to electron-volts,
\[
\phi
=
\frac{6.63\times10^{-19}}
{1.60\times10^{-19}}
\approx4.14\ \mathrm{eV}.
\]

Therefore
\[
\boxed{\phi\approx4.14\ \mathrm{eV}}.
\]

\subsection{(ii) Photoelectron current}

The incident laser has power
\[
P=100\ \mathrm W
\]
and wavelength
\[
\lambda=\SI{200}{nm}.
\]

The energy of each incident photon is
\[
E_\gamma=\frac{hc}{\lambda}.
\]

Hence
\[
E_\gamma
=
\frac{(6.63\times10^{-34})(2.998\times10^8)}
{200\times10^{-9}}
\approx9.94\times10^{-19}\ \mathrm J
\]
or
\[
E_\gamma\approx6.21\ \mathrm{eV}.
\]

The maximum kinetic energy of an emitted electron is therefore
\[
K_{\max}
=
E_\gamma-\phi
\approx6.21-4.14
=2.07\ \mathrm{eV}.
\]

In joules,
\[
K_{\max}\approx3.31\times10^{-19}\ \mathrm J.
\]

The problem states that the photoelectron current constitutes $10\%$ of the
incident power, so the power carried by the emitted electrons is
\[
P_e=0.10(100)=10\ \mathrm W.
\]

Thus the electron emission rate is
\[
\dot N_e=\frac{P_e}{K_{\max}}
\approx
\frac{10}{3.31\times10^{-19}}
\approx3.02\times10^{19}\ \mathrm{s^{-1}}.
\]

Assuming each electron contributes one elementary charge,
\[
I=e\dot N_e.
\]

Therefore
\[
I
=
(1.60\times10^{-19})(3.02\times10^{19})
\approx4.83\ \mathrm A.
\]

Hence
\[
\boxed{I\approx4.8\ \mathrm A}.
\]

% ============================================================
\section{Question 3 --- de Broglie wavelength of a water molecule}

The problem gives the molecular kinetic energy as
\[
E=\frac32 k_BT.
\]

For water,
\[
\mathrm{H_2O}
=
{}^1_1\mathrm H_2+{}^{16}_8\mathrm O.
\]

Ignoring electron masses and binding-energy mass defects, the mass is
\[
m
=
2m_p+8m_p+8m_n.
\]

Thus
\[
m
=
10m_p+8m_n.
\]

Using
\[
m_p=1.6726\times10^{-27}\ \mathrm{kg},
\qquad
m_n=1.6749\times10^{-27}\ \mathrm{kg},
\]
gives
\[
m\approx3.0125\times10^{-26}\ \mathrm{kg}.
\]

At
\[
T=300\ \mathrm K,
\]
the kinetic energy is
\[
E
=
\frac32(1.38\times10^{-23})(300)
\approx6.21\times10^{-21}\ \mathrm J.
\]

Using
\[
E=\frac{p^2}{2m},
\]
the momentum is
\[
p=\sqrt{2mE}.
\]

The de Broglie relation gives
\[
\lambda=\frac hp,
\]
so
\[
\lambda
=
\frac{h}{\sqrt{2mE}}
=
\frac{h}
{\sqrt{2m(3k_BT/2)}}.
\]

Therefore
\[
\lambda
=
\frac{6.63\times10^{-34}}
{\sqrt{
2(3.0125\times10^{-26})
(6.21\times10^{-21})
}}
\approx3.43\times10^{-11}\ \mathrm m.
\]

Thus
\[
\boxed{\lambda\approx3.43\times10^{-11}\ \mathrm m}
\]
or
\[
\boxed{\lambda\approx0.0343\ \mathrm{nm}}.
\]

This is much smaller than the size of an individual water molecule, illustrating why
the wave nature of matter is difficult to observe for macroscopic objects.

% ============================================================
\section{Question 4 --- Relativistic electrons}

The electron is accelerated until
\[
v=\frac c2.
\]

\subsection{(i) Accelerating voltage}

\subsubsection*{Classical calculation}

Classically,
\[
K=\frac12m_ev^2.
\]

Since an electron accelerated through a potential difference $V$ gains energy $eV$,
\[
eV=\frac12m_ev^2.
\]

Setting $v=c/2$,
\[
eV
=
\frac12m_e\frac{c^2}{4}
=
\frac{m_ec^2}{8}.
\]

Therefore
\[
V=\frac{m_ec^2}{8e}.
\]

Substitution gives
\[
V
\approx6.40\times10^4\ \mathrm V.
\]

Hence
\[
\boxed{V_{\rm classical}\approx64\ \mathrm{kV}}.
\]

\subsubsection*{Relativistic calculation}

The relativistic kinetic energy is
\[
K=(\gamma-1)m_ec^2,
\]
where
\[
\gamma=
\frac{1}{\sqrt{1-v^2/c^2}}.
\]

For $v=c/2$,
\[
\gamma
=
\frac{1}{\sqrt{1-\frac14}}
=
\frac{2}{\sqrt3}
\approx1.155.
\]

Therefore
\[
eV=(\gamma-1)m_ec^2,
\]
so
\[
V
=
\left(\frac{2}{\sqrt3}-1\right)
\frac{m_ec^2}{e}.
\]

Numerically,
\[
\boxed{V_{\rm relativistic}\approx79.2\ \mathrm{kV}}.
\]

The relativistic result is noticeably larger because the classical expression
underestimates the kinetic energy at this speed.

\subsection{(ii) de Broglie wavelength}

\subsubsection*{Classical momentum}

Classically,
\[
p_{\rm cl}=m_ev.
\]

Thus
\[
\lambda_{\rm cl}
=
\frac{h}{m_ev}.
\]

For $v=c/2$,
\[
\lambda_{\rm cl}
=
\frac{6.63\times10^{-34}}
{(9.109\times10^{-31})(c/2)}
\approx4.86\times10^{-12}\ \mathrm m.
\]

Therefore
\[
\boxed{\lambda_{\rm cl}\approx4.86\times10^{-12}\ \mathrm m}.
\]

\subsubsection*{Relativistic momentum}

The relativistic momentum is
\[
p=\gamma m_ev.
\]

Hence
\[
\lambda_{\rm rel}
=
\frac{h}{\gamma m_ev}.
\]

Using $\gamma=2/\sqrt3$,
\[
\lambda_{\rm rel}
\approx4.20\times10^{-12}\ \mathrm m.
\]

Thus
\[
\boxed{\lambda_{\rm rel}\approx4.20\times10^{-12}\ \mathrm m}.
\]

The relativistic wavelength is shorter because the relativistic momentum is larger
than the classical value at the same velocity.

\subsection{(iii) Atomic lattice spacing}

The electron beam is diffracted through an angle
\[
\theta=2^\circ.
\]

For first-order constructive interference,
\[
d\sin\theta=n\lambda,
\]
with
\[
n=1.
\]

Thus
\[
d=\frac{\lambda}{\sin2^\circ}.
\]

Using the more accurate relativistic wavelength,
\[
d
=
\frac{4.20\times10^{-12}}{\sin2^\circ}
\approx1.20\times10^{-10}\ \mathrm m.
\]

Therefore
\[
\boxed{d\approx1.2\times10^{-10}\ \mathrm m}.
\]

This is of the expected order of magnitude for an atomic lattice spacing.

% ============================================================
\section{Question 5 --- Ionisation energies and the Bohr model}

For a hydrogenic atom with nuclear charge $+Ze$ and one electron, the Bohr model
gives
\[
E_n=-\frac{13.6Z^2}{n^2}\ \mathrm{eV}.
\]

Ionisation from the ground state therefore requires
\[
E_{\rm ion}=13.6Z^2\ \mathrm{eV}.
\]

For a photon of energy $E$,
\[
E=\frac{hc}{\lambda}.
\]

Using $hc\approx1242\ \mathrm{eV\,nm}$,
\[
\lambda_{\max}
=
\frac{1242}{E_{\rm ion}}\ \mathrm{nm}.
\]

\subsection{(i) Maximum ionising wavelength}

\subsubsection*{(a) Hydrogen}

For $Z=1$,
\[
E_{\rm ion}=13.6\ \mathrm{eV}.
\]

Therefore
\[
\lambda_{\max}
=
\frac{1242}{13.6}
\approx91.3\ \mathrm{nm}.
\]

Thus
\[
\boxed{\lambda_{\max}(\mathrm H)\approx91.3\ \mathrm{nm}}.
\]

\subsubsection*{(b) Helium}

For the one-electron hydrogenic approximation with $Z=2$,
\[
E_{\rm ion}=13.6(2)^2
=54.4\ \mathrm{eV}.
\]

Hence
\[
\lambda_{\max}
=
\frac{1242}{54.4}
\approx22.8\ \mathrm{nm}.
\]

Therefore
\[
\boxed{\lambda_{\max}(\mathrm{He})\approx22.8\ \mathrm{nm}}.
\]

This is specifically the hydrogenic approximation; a real neutral helium atom has
two electrons, so the assumptions of the model do not apply directly.

\subsubsection*{(c) Radon}

For $Z=86$,
\[
E_{\rm ion}
=
13.6(86)^2
\approx1.006\times10^5\ \mathrm{eV}.
\]

Thus
\[
\lambda_{\max}
=
\frac{1242}{1.006\times10^5}
\approx1.24\times10^{-2}\ \mathrm{nm}.
\]

Therefore
\[
\boxed{\lambda_{\max}(\mathrm{Rn})\approx0.0124\ \mathrm{nm}}.
\]

Again, this is the formal hydrogenic prediction. For multi-electron atoms the
electron--electron repulsion and screening make the one-electron Bohr model
increasingly inappropriate.

\subsection{(ii)(a) Comparison with measured ionisation energies}

The measured molar ionisation energies are

\[
\begin{array}{lll}
\mathrm H &: &1312.0\ \mathrm{kJ\,mol^{-1}},\\
\mathrm{He}&: &2372.3\ \mathrm{kJ\,mol^{-1}},\\
\mathrm{Rn}&: &1037\ \mathrm{kJ\,mol^{-1}}.
\end{array}
\]

The hydrogen value is close to the Bohr prediction because hydrogen is genuinely
a one-electron atom. For helium and radon, however, the simple $13.6Z^2$ prediction
becomes increasingly poor.

The fundamental issue is that the Bohr expression used above assumes a
\emph{hydrogenic, one-electron atom}: the electron is moving in the Coulomb field of
a nucleus with charge $+Ze$, without accounting for the other electrons.

A measured ionisation energy for a neutral atom instead means removing one electron
from the actual multi-electron atom. Electron--electron repulsion and shielding
therefore have a major effect.

For example, the formal Bohr prediction for radon is
\[
\sim100.6\ \mathrm{keV},
\]
whereas its measured first ionisation energy is only about
\[
10.8\ \mathrm{eV}.
\]

This enormous discrepancy demonstrates the failure of applying the hydrogenic
formula directly to a many-electron atom.

\subsection{(ii)(b) Photon wavelengths for the measured ionisation energies}

The energy per atom corresponding to a molar ionisation energy $E_{\rm mol}$ is
\[
E=\frac{E_{\rm mol}}{N_A}.
\]

Equivalently, converting directly to electron-volts,
\[
E_{\rm eV}
=
\frac{E_{\rm mol}\times10^3}
{N_Ae}.
\]

The wavelength is then
\[
\lambda=\frac{hc}{E}.
\]

\subsubsection*{Hydrogen}

\[
E
=
\frac{1312.0\times10^3}
{6.022140857\times10^{23}}
\approx2.18\times10^{-18}\ \mathrm J,
\]
giving
\[
\boxed{\lambda_{\mathrm H}\approx91.2\ \mathrm{nm}}.
\]

\subsubsection*{Helium}

\[
E
=
\frac{2372.3\times10^3}
{6.022140857\times10^{23}},
\]
giving
\[
\boxed{\lambda_{\mathrm{He}}\approx50.5\ \mathrm{nm}}.
\]

\subsubsection*{Radon}

\[
E
=
\frac{1037\times10^3}
{6.022140857\times10^{23}},
\]
giving
\[
\boxed{\lambda_{\mathrm{Rn}}\approx115\ \mathrm{nm}}.
\]

The measured wavelengths therefore differ dramatically from the hydrogenic
prediction for the heavy, many-electron atom.

% ============================================================
\section{Question 6 --- FAST and photons from Andromeda}

FAST has diameter
\[
D=500\ \mathrm m,
\]
so its collecting area is
\[
A=\pi\left(\frac D2\right)^2
=\pi(250)^2.
\]

The distance to Andromeda is
\[
r=2\times10^{22}\ \mathrm m,
\]
and the power radiated at
\[
f=1420\ \mathrm{MHz}
=1.420\times10^9\ \mathrm{Hz}
\]
is
\[
P_{\rm A}=8\times10^{27}\ \mathrm W.
\]

\subsection{(i) Photons received by FAST}

Assuming isotropic radiation, the power flux at FAST is
\[
I_{\rm A}
=
\frac{P_{\rm A}}{4\pi r^2}.
\]

The power intercepted by FAST is
\[
P_{\rm FAST}
=
I_{\rm A}A.
\]

Combining these,
\[
P_{\rm FAST}
=
\frac{P_{\rm A}}{4\pi r^2}
\pi\left(\frac D2\right)^2
=
\frac{P_{\rm A}}{4}
\left(\frac{D}{2r}\right)^2.
\]

Numerically,
\[
P_{\rm FAST}\approx3.13\times10^{-13}\ \mathrm W.
\]

Each photon has energy
\[
E_\gamma=hf.
\]

Thus
\[
E_\gamma
=
(6.63\times10^{-34})(1.420\times10^9)
\approx9.41\times10^{-25}\ \mathrm J.
\]

Therefore the number of photons received per second is
\[
\dot N_{\rm FAST}
=
\frac{P_{\rm FAST}}{hf}
\approx3.3\times10^{11}\ \mathrm{s^{-1}}.
\]

Hence
\[
\boxed{\dot N_{\rm FAST}\approx3.3\times10^{11}\ \mathrm{photons\,s^{-1}}}.
\]

\subsection{(ii) DSLR exposure}

The DSLR sensor has dimensions
\[
24\,\mathrm{mm}\times36\,\mathrm{mm},
\]
so its area is
\[
A_{\rm DSLR}
=
(24\times10^{-3})(36\times10^{-3})
=
8.64\times10^{-4}\ \mathrm{m^2}.
\]

With solar intensity
\[
I_\odot=100\ \mathrm{W\,m^{-2}},
\]
the incident power is
\[
P_{\rm DSLR}
=
I_\odot A_{\rm DSLR}
=
0.0864\ \mathrm W.
\]

For a shutter time of
\[
t=\frac1{500}\ \mathrm s,
\]
the received energy is
\[
E_{\rm DSLR}
=
P_{\rm DSLR}t
=
1.728\times10^{-4}\ \mathrm J.
\]

For sunlight of wavelength
\[
\lambda=500\ \mathrm{nm},
\]
the photon energy is
\[
E_\gamma
=
\frac{hc}{\lambda}
\approx3.98\times10^{-19}\ \mathrm J.
\]

Thus the number of photons in the exposure is
\[
N_{\rm DSLR}
=
\frac{E_{\rm DSLR}}{E_\gamma}
\approx4.35\times10^{14}.
\]

Therefore
\[
\boxed{N_{\rm DSLR}\approx4.35\times10^{14}\ \mathrm{photons}}.
\]

\subsection{(iii) Equivalent FAST exposure}

FAST receives approximately
\[
3.3\times10^{11}\ \mathrm{photons\,s^{-1}}.
\]

The time needed to receive the same number of photons as the DSLR exposure is
\[
t_{\rm FAST}
=
\frac{N_{\rm DSLR}}{\dot N_{\rm FAST}}.
\]

Thus
\[
t_{\rm FAST}
\approx
\frac{4.35\times10^{14}}
{3.3\times10^{11}}
\approx1.32\times10^3\ \mathrm s.
\]

Therefore
\[
\boxed{t_{\rm FAST}\approx22\ \mathrm{minutes}}.
\]

The comparison is striking: although Andromeda is enormously luminous, its huge
distance makes the photon flux arriving at Earth extremely small.

% ============================================================
\section{Question 7 --- Solar sail}

The incident light has intensity
\[
I\ \mathrm{W\,m^{-2}}
\]
and the sail has area $A$.

\subsection{(i) Impulse delivered per second}

For a photon of wavelength $\lambda$, the de Broglie relation gives
\[
p=\frac{h}{\lambda}.
\]

When a photon is perfectly reflected, its momentum reverses direction. Therefore
the change in momentum of the photon, and hence the impulse delivered to the sail,
is
\[
\Delta p=2p=\frac{2h}{\lambda}.
\]

The energy of one photon is
\[
E_\gamma=hf=\frac{hc}{\lambda}.
\]

The number of photons arriving at the sail per second is therefore
\[
\dot N
=
\frac{IA}{hc/\lambda}
=
\frac{IA\lambda}{hc}.
\]

The impulse delivered per second is
\[
F
=
\dot N\Delta p.
\]

Hence
\[
F
=
\frac{IA\lambda}{hc}
\frac{2h}{\lambda}
=
\boxed{\frac{2IA}{c}}.
\]

This is the radiation-pressure force for a perfectly reflecting sail.

\subsection{(ii) Acceleration}

Newton's second law gives
\[
F=ma.
\]

Therefore
\[
ma=\frac{2IA}{c},
\]
so
\[
\boxed{
a=\frac{2IA}{mc}
}.
\]

The corresponding radiation pressure is
\[
P_{\rm rad}=\frac FA=\boxed{\frac{2I}{c}}.
\]

\subsection{(iii) Numerical example}

The spacecraft has mass
\[
m=10\ \mathrm{tonnes}=10^4\ \mathrm{kg},
\]
and the sail area is
\[
A=100\ \mathrm{km^2}=10^8\ \mathrm{m^2}.
\]

The intensity is
\[
I=1\ \mathrm{kW\,m^{-2}}=10^3\ \mathrm{W\,m^{-2}}.
\]

Thus
\[
a
=
\frac{2(10^3)(10^8)}
{(10^4)(2.998\times10^8)}
\approx6.67\times10^{-2}\ \mathrm{m\,s^{-2}}.
\]

Therefore
\[
\boxed{a\approx0.0667\ \mathrm{m\,s^{-2}}}.
\]

The target velocity is
\[
v=\frac{c}{100}.
\]

Assuming constant acceleration,
\[
t=\frac va
=
\frac{c/100}{a}.
\]

Hence
\[
t
\approx
\frac{2.998\times10^8}
{100(0.0667)}
\approx4.49\times10^7\ \mathrm s.
\]

Converting to years,
\[
t\approx1.42\ \mathrm{yr}.
\]

Thus
\[
\boxed{t\approx4.49\times10^7\ \mathrm s\approx1.4\ \mathrm{years}}.
\]

The problem explicitly notes that this constant-acceleration treatment eventually
breaks down as relativistic effects become important.

% ============================================================
\section{Question 8 --- Stellar radiation and mass loss}

\subsection{(i) Solar luminosity and mass loss}

The Earth receives a maximum solar intensity of
\[
I=1361\ \mathrm{W\,m^{-2}}.
\]

At one astronomical unit,
\[
r=1.496\times10^{11}\ \mathrm m.
\]

Assuming the Sun radiates isotropically, its luminosity is
\[
L_\odot=4\pi r^2I.
\]

Thus
\[
L_\odot
=
4\pi(1.496\times10^{11})^2(1361)
\approx3.83\times10^{26}\ \mathrm W.
\]

Hence
\[
\boxed{L_\odot\approx3.83\times10^{26}\ \mathrm W}.
\]

Over
\[
t=10^{10}\ \mathrm{yr},
\]
the radiated energy is
\[
E=L_\odot t.
\]

Using
\[
1\ \mathrm{yr}\approx365\times24\times3600\ \mathrm s,
\]
we obtain
\[
E\approx1.21\times10^{44}\ \mathrm J.
\]

Einstein's mass-energy relation gives
\[
E=\Delta mc^2,
\]
so
\[
\Delta m
=
\frac{E}{c^2}
\approx1.34\times10^{27}\ \mathrm{kg}.
\]

The Sun's mass is
\[
M_\odot=1.989\times10^{30}\ \mathrm{kg},
\]
so the percentage mass loss is
\[
\frac{\Delta m}{M_\odot}\times100
\approx0.067\%.
\]

Therefore
\[
\boxed{\text{solar mass loss over this period}\approx0.07\%}.
\]

\subsection{(ii) Sirius}

Sirius has
\[
T_S=9940\ \mathrm K,
\qquad
R_S=1.711R_\odot,
\qquad
M_S=2.02M_\odot.
\]

\subsubsection*{(a) Radiation received at Earth's orbit}

For the same distance from the star, the received intensity is proportional to
\[
T^4R^2.
\]

Therefore
\[
\frac{I_S}{I_\odot}
=
\left(\frac{T_S}{T_\odot}\right)^4
\left(\frac{R_S}{R_\odot}\right)^2.
\]

Thus
\[
I_S
=
1361
\left(\frac{9940}{5778}\right)^4
(1.711)^2.
\]

This gives
\[
\boxed{I_S\approx3.49\times10^4\ \mathrm{W\,m^{-2}}}
\]
or approximately
\[
\boxed{34.9\ \mathrm{kW\,m^{-2}}}.
\]
This is also obtained directly from the luminosity of Sirius divided by
$4\pi(1\,\mathrm{AU})^2$.

\subsubsection*{(b) Luminosity of Sirius}

The Stefan--Boltzmann law gives
\[
L_S=4\pi R_S^2\sigma T_S^4.
\]

Relative to the Sun,
\[
\frac{L_S}{L_\odot}
=
\left(\frac{R_S}{R_\odot}\right)^2
\left(\frac{T_S}{T_\odot}\right)^4.
\]

Therefore
\[
L_S
=
3.83\times10^{26}
(1.711)^2
\left(\frac{9940}{5778}\right)^4.
\]

Hence
\[
\boxed{L_S\approx9.82\times10^{27}\ \mathrm{J\,s^{-1}}}
\]
or
\[
\boxed{L_S\approx9.82\times10^{27}\ \mathrm W}.
\]

\subsubsection*{(c) Lifetime}

The problem asks us to assume that Sirius loses the same \emph{fraction} of its
mass over its lifetime as the Sun does during its main-sequence phase. Thus the
fractional mass loss is approximately
\[
f=0.07\%=7.0\times10^{-4}.
\]
For Sirius, the mass converted to radiated energy is therefore
\[
\Delta M_S=fM_S=f(2.02M_\odot).
\]
Using $E=\Delta Mc^2$ and the Sirius luminosity from part (b),
\[
t_S=\frac{\Delta M_Sc^2}{L_S}.
\]
Hence
\[
t_S
=\frac{(7.0\times10^{-4})(2.02)(1.989\times10^{30})(2.998\times10^8)^2}
{9.82\times10^{27}}
\approx2.57\times10^{16}\,\mathrm{s}.
\]
Converting to years,
\[
\boxed{t_S\approx8.2\times10^8\,\mathrm{yr}\approx0.82\,\mathrm{billion\ years}}.
\]
This is not a realistic stellar lifetime calculation: real stellar evolution changes the
luminosity, composition and mass-loss rate over time. The result follows from the
specific constant-fractional-mass-loss assumption in the question.

\subsection{(iii) DogStar heating elements}

A $10\,\mathrm{m^2}$ panel at one astronomical unit from Sirius receives power
\[
P=I_S A.
\]

Using the irradiance found in part (a),
\[
I_S\approx3.49\times10^4\,\mathrm{W\,m^{-2}},
\]
we find
\[
P
\approx
(3.49\times10^4)(10)
\approx3.49\times10^5\,\mathrm W.
\]
Thus the heating elements are idealised as having an available power of
\[
\boxed{P\approx3.49\times10^5\,\mathrm W}.
\]

To boil one litre of water from $20^\circ\mathrm C$ to $100^\circ\mathrm C$,
\[
Q=mc\Delta T.
\]

With
\[
m\approx1\ \mathrm{kg},
\qquad
c=4181\ \mathrm{J\,kg^{-1}\,K^{-1}},
\qquad
\Delta T=80\ \mathrm K,
\]
we obtain
\[
Q
=
(1)(4181)(80)
=
3.345\times10^5\ \mathrm J.
\]

Using the supplied solution's power estimate,
\[
P\approx3.49\times10^5\ \mathrm W,
\]
the heating time is
\[
t=\frac QP
\approx
\frac{3.345\times10^5}{3.49\times10^5}
\approx0.96\ \mathrm s.
\]

Therefore
\[
\boxed{t\approx1\ \mathrm s}.
\]

This ignores heat losses, the finite heat capacity of the container, evaporation,
and the energy required to maintain the water at its boiling point; it is therefore
an idealised minimum heating time.

% ============================================================
\section*{Summary of numerical results}

\begin{center}
\small
\begin{tabular}{p{0.17\textwidth}p{0.70\textwidth}}
\toprule
\textbf{Part} & \textbf{Result}\\
\midrule
1(i) & $\lambda\approx3.26\,\mathrm{cm}$, microwave\\
1(ii) & $\Delta E\approx3.81\times10^{-5}\,\mathrm{eV}$\\
2(i) & $\phi\approx4.14\,\mathrm{eV}$\\
2(ii) & $I\approx4.8\,\mathrm A$\\
3 & $\lambda_{\rm H_2O}\approx3.43\times10^{-11}\,\mathrm m$\\
4(i) & $V_{\rm cl}\approx64\,\mathrm{kV}$; $V_{\rm rel}\approx79.2\,\mathrm{kV}$\\
4(ii) & $\lambda_{\rm cl}\approx4.86\,\mathrm{pm}$; $\lambda_{\rm rel}\approx4.20\,\mathrm{pm}$\\
4(iii) & $d\approx1.2\times10^{-10}\,\mathrm m$\\
5(i) & $\lambda_{\rm H}\approx91.3\,\mathrm{nm}$; $\lambda_{\rm He}\approx22.8\,\mathrm{nm}$; $\lambda_{\rm Rn}\approx0.0124\,\mathrm{nm}$\\
5(ii) & Measured wavelengths: H $\approx91.2\,\mathrm{nm}$, He $\approx50.5\,\mathrm{nm}$, Rn $\approx115\,\mathrm{nm}$\\
6 & FAST $\approx3.3\times10^{11}$ photons/s; equivalent DSLR exposure $\approx22$ min\\
7 & $a\approx0.0667\,\mathrm{m\,s^{-2}}$; $t\approx1.4$ yr to reach $c/100$\\
8(i) & $L_\odot\approx3.83\times10^{26}\,\mathrm W$; mass loss $\approx0.07\%$\\
8(ii)(a) & Sirius irradiance $\approx3.49\times10^4\,\mathrm{W\,m^{-2}}$\\
8(ii)(b) & $L_S\approx9.82\times10^{27}\,\mathrm W$\\
8(ii)(c) & Simplified lifetime $\approx8.2\times10^8$ yr\\
8(iii) & Idealised boiling time $\approx1$ s\\
\bottomrule
\end{tabular}
\end{center}

\end{document}
